\section{摘要}

\subsection{研究背景与动机}

乳腺癌作为全球女性发病率最高的恶性肿瘤，其分子亚型的精准识别对个体化治疗决策具有至关重要的临床意义。PAM50分类系统将乳腺癌划分为Luminal A（LumA）、Luminal B（LumB）、HER2富集型和Basal-like四个内在亚型，不同亚型在治疗响应和预后方面存在显著差异。传统基于单一组学数据的分类方法难以充分刻画肿瘤的多层次分子异质性，而多组学融合策略通过整合转录组、表观遗传组等多维度分子信息，为提升亚型分类精度提供了新的技术路径。然而，多组学融合方法在小样本条件下的性能表现尚缺乏系统的实证研究，特别是复杂融合方法（如Stacking）是否必然优于简单融合方法（如Concat）这一核心问题，亟待通过严格的实验设计加以回答。

\subsection{研究方法与技术路线}

本研究基于TCGA-BRCA数据库的RNA-seq转录组、DNA甲基化以及PAM50临床标签数据，围绕"单癌种内部分子亚型分类"这一核心任务，系统比较了RNA单模态基线、早期融合（Concat）、潜在因子融合（MOFA）与堆叠集成（Stacking）四种多组学融合策略的分类性能。为确保实验结果的科学可信度，本研究构建了严格的防泄露评估框架：所有特征筛选、数据标准化、MOFA因子拟合及Stacking元特征构造均限定在训练折内完成，测试折仅用于独立推断，从而消除了因数据泄露导致的性能高估风险。主实验采用重复分层5折交叉验证（5 folds $\times$ 10 repeats，共50个测试折）作为统一评估口径，以Accuracy、Balanced Accuracy和Macro-F1作为核心评价指标，并报告95\%置信区间以量化估计不确定性。此外，本研究还系统开展了特征维度敏感性（100--1500维）、CV参数敏感性（3/5/10折、3--15次重复）、MOFA因子数敏感性（10--25个因子）以及消融实验（去除RNA/去除甲基化/无特征选择/类别平衡）等补充实验，共计41组实验配置，为方法选择提供了全面的量化依据。

\subsection{主要实验结果}

实验结果表明，在有效样本量约150例、三分类（LumA/LumB/Basal）的任务设定下，早期融合策略Concat-SVM取得了最优的综合性能，其Accuracy达到90.64\%（95\% CI: [87.75\%, 93.53\%]），Balanced Accuracy为90.56\%，Macro-F1为90.03\%。MOFA（20个潜在因子）表现次之，Accuracy为90.00\%，与Concat-SVM的差距不足1个百分点。RNA单模态基线已达到89.78\%的准确率，表明转录组数据本身蕴含丰富的亚型判别信息。相比之下，Stacking集成策略表现令人意外：被设计为"SOTA"的XGB+XGB配置仅取得78.66\%的Macro-F1，而所有Stacking变体的性能均显著低于简单融合基线，提示在小样本高维条件下增加融合层级反而引入过拟合风险。消融实验进一步揭示：去除甲基化后Concat性能下降约2.2个百分点，而去除RNA后性能骤降10.6个百分点，确认RNA数据是分类性能的主要贡献来源，甲基化数据则起到重要的互补稳定作用。

\subsection{核心结论与贡献}

本研究的核心结论包括：第一，在小样本条件下，简单融合方法（Concat-SVM）优于复杂融合方法（Stacking），挑战了"越复杂越好"的直觉假设；第二，RNA转录组数据是乳腺癌分子亚型分类的主导信息源，甲基化数据作为补充可带来有限但稳定的性能提升；第三，严格的防泄露评估框架是保证方法比较公平性的关键，任何信息泄露都可能导致性能高估和错误结论。本研究的贡献在于建立了可审计、可复现的多组学融合评估框架，并首次在同一数据集上系统揭示了"简单方法优于复杂方法"的反直觉现象，为乳腺癌分子亚型分类的临床转化提供了方法论参考。同时，本研究也明确了当前局限性（样本量有限、单一数据集、HER2类别样本稀少），并提出了外部验证、深度学习方法探索、生物学机制阐释等后续研究方向。
